% !TeX root = ../papers/nc-p3.tex
\section{整な階数二曲線の両端の次数と核像}
\label{sec:cofactor-extreme-degrees}
\begin{theorem}\label{thm:cofactor-exclude-seventeen}
定理\ref{thm:cofactor-integral-octic}の整な階数二曲線の次数十七は生じない。
\end{theorem}
\begin{proof}
$d=17$ とすると前定理のClifford不等式で $A=H^0(\mathcal O_C(4))$ の次元はちょうど三十五、
$B=H^0(\omega_C(-4))\ne0$、$g=p_a(C)=34+\dim B$ である。
$\dim B\ge2$ と仮定する。滑らかな一点で付値の異なる $B$ の基底を選び、最大の二元を $b_0,b_1$ とする。
残りを付値の降順に積の包に加えると、$A$ の最小付値との積が既存の包のどの元よりも
小さい付値を持つので、次元が毎回一以上増える。
積の全体は $g$ 次元に入るから $\dim(A+fA)\le36$, $f=b_1/b_0$。$f$ は非定数なので等号である。
分数体で $A$ を自明化し、$A_j=\{a\in A:f^ia\in A,\ 0\le i\le j\}$ とする。
$f^jA_{j-1}\subset fA$ より $A_{j-1}/A_j$ の次元は高々一。
非零の間に等号 $A_j=A_{j-1}$ があれば $f A_{j-1}\subset A_{j-1}$ となり不可能である。
ゆえにある非零元の $1,f,\ldots,f^{34}$ 倍が $A$ の基底となる。
その射影像は次数三十四の有理正規曲線であるが、$A$ は次数六十八の非常に豊富な線束を与えるので矛盾。
従って $\dim B=1$, $g=35$。$\mathcal O_C(4)\hookrightarrow\omega_C$ は同じEuler数を持つ
階数一層の包含ゆえ同型であり、$C$ はGorensteinである。
正則三次元局所環内の余次元二Gorensteinイデアルは完全交叉である
（長さ二の極小分解の最終階数が一で、イデアルの生成数は二）。

命題\ref{prop:cofactor-singular-length}で $2m\ge102-68+5=39$ なので二十個の印は全て $C$ に載る。
純外積層 $W$ はPlücker非退化性から $h^0(W)\ge6$。
定理\ref{thm:cofactor-curve-duality}と既約八次式の余因子が七以上独立なこと
（定理\ref{thm:seven-cofactor-trisecant-grassmannian}の核像の証明）から $h^1(W)\le3$。
従って $\chi(W)\ge3$ で、$a=80-6d-\chi(\mathcal O_C)-\chi(W)\le9$。
少なくとも十一個の印で $a_p=0$、すなわち余因子が局所的に $C$ のイデアルそのものを定める。
各印ではこのイデアルは極大イデアルの二乗に入り、二生成の完全交叉なので曲線の重複度 $e_p\ge4$。
実際一般平面で二つの位数二以上の式を切れば局所交点数は四以上である。
正規化の局所余長はこの各点で四以上になる。一般線形パラメータ $t$ に対し正規化半局所環
$\widetilde R$ では $\mathfrak m\widetilde R=t\widetilde R$、$\ell(\widetilde R/t^2\widetilde R)=2e_p$。
そこへの元の環の像は $R/\mathfrak m^2$ を経るので長さ四以下、よって余長は $2e_p-4\ge4$。
十一点で四十四となり $g=35$ に反する。
\end{proof}

\begin{theorem}\label{thm:cofactor-sectional-kernel}
残る次数 $d=18,19,20,21$ とし、$D$ を既約八次行列式曲面、$K$ を階数三での核像の閉包とする。
一般平面八次切断の正規化で核線束の次数は $24-d$、種数は $21-d$ 以下である。
$K$ が曲線ならその次数は $24-d$ である。$K$ が曲面で主成分の次数が $n$ なら
\[
                 8n\deg K\le(24-d)^2 .
\]
特に $d=21$ は必ず七余因子のBezout標準形である。$d=20$ で曲面像となる場合は
$K$ は既約二次曲面であり、核写像は双有理である。
\end{theorem}
\begin{proof}
一般入力平面切断の正規化上で六次余因子列は核点の平方を共通因子倍したものである。
階数三の点では共通零がない。階数二の $d$ 点は曲線の一般点で、余因子イデアルは被約曲線のイデアル、
平面では各点の極大イデアルである。平面八次のその点での重複度はちょうど二である
（曲線の一般点での二次初項と横断平面を使う）。
従って平面八次の全正規化枝での極大イデアルの位数和は二であり、余因子の共通零の総次数は $2d$。
単純節点を仮定せず $2\deg L=48-2d$ を得る。種数上界は平面八次の算術種数二十一による。
核グラフの滑らかな解消で、入力と出力の超平面引戻しを $H,L$ と書けば $H^2=8$, $HL=24-d$。
実際一般入力平面切断は例外曲線を成分に持たず、種数と次数はその正規化で計算できる。
核像が曲線なら定理\ref{thm:kernel-curve-scroll}で一般繊維は一本の入力直線、
従って $HL=\deg K$。曲面像なら $L^2=n\deg K$ でHodge指数定理が不等式を与える。
$K$ は非退化なので曲面なら次数二以上である。$d=21$ では不等式と両立せず、
曲線像は次数三だから有理正規三次曲線。定理\ref{thm:seven-cofactor-trisecant-grassmannian}の
Bezout表示と入力・出力の基底変更判定が適用される。$d=20$ の曲面像では数値から $n=1,\deg K=2$。
\end{proof}
